\begin{table}[!ht]
\centering
\begin{threeparttable}
\caption{Overfitting Diagnostic: Train versus Test RMSE}
\label{tab:overfitting}

\begin{tabular}{llrrr}
\toprule
Model & Horizon & Train RMSE & Test RMSE & Gap\\
\midrule
OLS & 1d & 0.01607 & 0.01644 & -0.00037\\
OLS & 3d & 0.02543 & 0.02651 & -0.00108\\
OLS & 5d & 0.03164 & 0.03322 & -0.00158\\
Random Forest & 1d & 0.01454 & 0.01678 & -0.00224\\
Random Forest & 3d & 0.02298 & 0.02716 & -0.00418\\
Random Forest & 5d & 0.02757 & 0.03360 & -0.00603\\
XGBoost & 1d & 0.01610 & 0.01631 & -0.00021\\
XGBoost & 3d & 0.02533 & 0.02649 & -0.00116\\
XGBoost & 5d & 0.03049 & 0.03262 & -0.00213\\
SVM & 1d & 0.01613 & 0.01634 & -0.00021\\
SVM & 3d & 0.02583 & 0.02621 & -0.00038\\
SVM & 5d & 0.03245 & 0.03253 & -0.00008\\
ANN & 1d & 0.01619 & 0.01628 & -0.00009\\
ANN & 3d & 0.02562 & 0.02634 & -0.00072\\
ANN & 5d & 0.03203 & 0.03275 & -0.00072\\
\bottomrule
\end{tabular}
\begin{tablenotes}
\footnotesize
\item \textit{Note.} Gap = Train RMSE minus Test RMSE. Positive values indicate the model fits training data more closely than it generalises to unseen observations. Large gaps relative to test RMSE suggest overfitting.
\end{tablenotes}
\end{threeparttable}
\end{table}
